% Iteration 6 change manifest as two structured cards: chg-1 extends the
% publish-state guard from deliverable files to script entrypoints, chg-2
% introduces the cross-step ExecutionRiskHintsMiddleware. Both ship together
% as the iteration-6 harness. Palette and layout follow case-trajectory.tex.
\begin{figure}[t]
    \centering
    \scriptsize
    \begin{tcbraster}[raster columns=2,
                      raster equal height,
                      raster column skip=4pt,
                      raster row skip=0pt,
                      raster left skip=0pt,
                      raster right skip=0pt,
                      boxrule=0.6pt,
                      left=4pt,right=4pt,top=3pt,bottom=3pt,
                      coltitle=white]
        \begin{tcolorbox}[colback=aheSky!18,colframe=aheNavy,colbacktitle=aheNavy,
                          title={\scriptsize\textbf{chg-1，迭代 6，commit \texttt{ff0cf3d}，层级：工具实现}}]
            \textbf{\textcolor{aheNavy}{涉及文件}} \\
            \rule{\linewidth}{0.3pt}\\[1pt]
            $\bullet$ \texttt{tools/shell\_tools/run\_shell\_command.py} \\
            $\bullet$ \texttt{tool\_descriptions/run\_shell\_command.tool.yaml} \\[3pt]
            \rule{\linewidth}{0.3pt}\\
            \textbf{\textcolor{aheNavy}{改动内容}} \\
            \rule{\linewidth}{0.3pt}\\[1pt]
            在迭代 5 已经覆盖的交付文件与根目录之上，把发布态目标的提取范围扩展到脚本入口以及最终检查中显式引用的文件。在评测器式的最终检查成功之后，守卫如今除了拦截删除与根目录重置这两类情形，还会拦截对受保护文件的重写以及对受保护生成脚本的重跑。 \\[3pt]
            \rule{\linewidth}{0.3pt}\\
            \textbf{\textcolor{aheNavy}{修复的失败模式}} \\
            \rule{\linewidth}{0.3pt}\\[1pt]
            智能体带着一个已收敛的生成脚本进入发布态，随后以``收拾整洁''为名重跑或重写该脚本，使已经验证过的输出失效；这正是 \texttt{mcmc-sampling-stan} 的 F4 步骤。 \\[3pt]
            \rule{\linewidth}{0.3pt}\\
            \textbf{\textcolor{aheNavy}{预测修复的任务}} \\
            \rule{\linewidth}{0.3pt}\\[1pt]
            \texttt{mcmc-sampling-stan}，以及诸如 \texttt{configure-git-webserver} 这类残余的``验证之后又改动''的情形。
        \end{tcolorbox}
        \begin{tcolorbox}[colback=aheTeal!12,colframe=aheBlue,colbacktitle=aheBlue,
                          title={\scriptsize\textbf{chg-2，迭代 6，commit \texttt{9651986}，层级：中间件}}]
            \textbf{\textcolor{aheBlue}{涉及文件}} \\
            \rule{\linewidth}{0.3pt}\\[1pt]
            $\bullet$ \texttt{workspace/code\_agent.yaml} \\
            $\bullet$ \texttt{workspace/middleware/\_\_init\_\_.py} \\
            $\bullet$ \texttt{workspace/middleware/execution\_risk\_hints.py} \\[3pt]
            \rule{\linewidth}{0.3pt}\\
            \textbf{\textcolor{aheBlue}{改动内容}} \\
            \rule{\linewidth}{0.3pt}\\[1pt]
            通过一个 \texttt{AfterToolHook} 注册了新的 \texttt{ExecutionRiskHintsMiddleware}：它扫描每一条 shell 命令及其结果，跨步骤累积轻量状态，并在实时历史匹配上七种风险模式之一时排入一条针对性的提醒：借助 \texttt{--help}、\texttt{py\_compile} 或仅检查是否存在的浅层验证；契约指明了外部接口，却只在 localhost 上做服务检查；用内联或自写的代理验证器代替指定的评测器；绕过官方封装直接调用底层模型 API；基准运行时没有显式的黄金答案或阈值比较器；在同一命令形态上反复出现长时间超时；反复重试却始终撞上同一错误特征。提醒会去重，并在每次 rollout 内设置数量上限。 \\[3pt]
            \rule{\linewidth}{0.3pt}\\
            \textbf{\textcolor{aheBlue}{修复的失败模式}} \\
            \rule{\linewidth}{0.3pt}\\[1pt]
            只有从实时命令历史中才会显现的跨步骤行为，仅提示词的规则来不及对其作出反应。 \\[3pt]
            \rule{\linewidth}{0.3pt}\\
            \textbf{\textcolor{aheBlue}{预测修复的任务}} \\
            \rule{\linewidth}{0.3pt}\\[1pt]
            6 个任务。例如：\texttt{caffe-cifar-10}、\texttt{sam-cell-seg}、\texttt{mteb-retrieve}、\texttt{dna-assembly}、\texttt{train-fasttext}。
        \end{tcolorbox}
    \end{tcbraster}
    \caption{作为迭代 6 的 harness 上线的两条变更清单条目。\texttt{chg-1} 把迭代 5 的发布态守卫从交付文件扩展到脚本入口，这正是保护 \texttt{mcmc-sampling-stan} 中 \texttt{analysis.R} 所缺的一环。\texttt{chg-2} 引入了本次运行中的第一个跨步骤组件，即在实时命令历史上监视七种风险模式的 \texttt{ExecutionRiskHintsMiddleware}。}
    \label{fig:manifest-middleware}
\end{figure}
